\section{Exact generator control on the finite graph and its strict mixed quotient}
\label{sec:pggt}

Keep every source, coordinate and full-unit factor of
PGS.1--41, PGB.1--26 and PGMB.1--10. In the original increasing
remainder frame of $E=\mathbb{C}[S]/(\chi_{h,k})$, retain
\[
 A=M_S,\quad e_0=[1],\quad
 \ell[P]=[S^{q-1}]\operatorname{rem}_{\chi}P,\quad
 R=e_0\ell,\quad T_t=A+tR-kI/2,\qquad t\in\mathbb{C}.
 \tag{PGGT.1}\label{eq:pggt-original-generator}
\]
The multiplication $A$ retains every original primary
nilpotent. The parameter $t$ is the original period parameter;
it is not the real metric interpolation parameter, the
holonomy coordinate, or one of the original $t_i$ variables.

There are three actual positive-form comparisons on this same
coefficient space. The first two fix the original packet and
finite cutoff $N\geq q-1$; the averaged comparison also retains
the common original circumference domain of PGB.8. The third
uses the calculated completed quotient, with its original
two weighted observation spaces from PGMB.1 and PGMB.6:
\[
\begin{array}{c|c|c}
 G&H&B=G^{-1}(H-G)\\ \hline
 G_N&\widehat G_N&G_N^{-1}\mathcal E_N^{\rm aug}\\
 \overline G_N&\widehat G_N&
 \overline G_N^{-1}(\mathcal E_N^{\rm aug}+\mathscr D_N^{\rm hol})\\
 K_{\rm low}&\Omega_k&K_{\rm low}^{-1}
                                  (\Omega_k-K_{\rm low})
\end{array}
\tag{PGGT.2}\label{eq:pggt-three-comparisons}
\]
All three rows contain the forms already calculated on their
original sources. The complete PGB.6 and PGB.12 injectivity
proofs make the first two $B$ strictly positive in their
displayed $G$ metrics. The PGM.5--9 proof makes the third
strictly positive for $k\geq2$ and exactly zero for $k=1$.
Directly $GB=H-G$ is Hermitian, so $B^{\dagger_G}=B$;
and $x^*GBx=x^*(H-G)x$ proves each asserted sign in that
same original metric. None of the three rows identifies a
different source or completion with its neighbour.

For each row write the entire $G$-orthogonal spectral
decomposition as $B=\sum_{i=1}^r b_iP_i$, where
$0\leq b_1<\cdots<b_r$, each $P_i$ retains its full
multiplicity, $\sum_iP_i=I$, and $P_iP_j=\delta_{ij}P_i$.
These statements follow by diagonalizing the Hermitian
matrix $G^{1/2}BG^{-1/2}$ and transporting its complete
eigenspaces by $G^{-1/2}$. The transport changes no form:
the projections satisfy $P_i^*G=GP_i$. Define
\[
 s_i=\sqrt{1+b_i},\quad
 \mathsf S=\sum_i s_iP_i,\qquad
 \mathsf S^{-1}=\sum_i s_i^{-1}P_i,\quad
 \alpha=1+b_1>0.
 \tag{PGGT.3}\label{eq:pggt-exact-isometry}
\]
Then $\mathsf S^*G\mathsf S=G(I+B)=H$, by multiplying
the displayed sums. Thus this specific $\mathsf S$ is an
isometry $(E,H)\to(E,G)$ with the displayed inverse.
It need not commute with the original generator. That
failure is computed next rather than omitted.

Put $X_H=T_t^{\dagger_H}+T_t$,
$X_G=T_t^{\dagger_G}+T_t$ and
$\epsilon_H(t)=\|X_H\|_H$,
$\epsilon_G(t)=\|X_G\|_G$. The full transported error is
\[
\begin{gathered}
 Z=[\mathsf S,T_t]\mathsf S^{-1}
   =\sum_{i,j}\frac{s_i-s_j}{s_j}P_iT_tP_j,\\
 \mathsf S X_H\mathsf S^{-1}-X_G
                           =Z+Z^{\dagger_G},\\
 [B,T_t]=[B,A]+t[B,R]
            =\sum_{i,j}(b_i-b_j)P_iT_tP_j.
\end{gathered}
\tag{PGGT.4}\label{eq:pggt-entire-error}
\]
For the first line multiply the complete sums in PGGT.3 on
both sides of $T_t$. For the second, $H=G\mathsf S^2$
gives $T_t^{\dagger_H}=\mathsf S^{-2}
T_t^{\dagger_G}\mathsf S^2$. Therefore the left side
is $\mathsf S^{-1}T_t^{\dagger_G}\mathsf S
+\mathsf ST_t\mathsf S^{-1}-T_t^{\dagger_G}-T_t$,
whose two ordered differences are respectively
$Z^{\dagger_G}$ and $Z$. The last line follows from
$[B,-kI/2]=0$ and the complete spectral multiplication.
This retains the original complex $t$; taking the adjoint
in the second line consequently retains $\bar t$ as well.

The exact $G$-operator norm of the Hermitian matrix in
PGGT.4 is one control with all of these cancellations. A
second directly evaluated control is
\[
\begin{gathered}
 |\epsilon_H(t)-\epsilon_G(t)|
 \leq\|Z+Z^{\dagger_G}\|_G
 \leq\frac{\|[B,A]+t[B,R]\|_G}{1+b_1}
 \leq\frac{b_r-b_1}{1+b_1}\|T_t\|_G.
\end{gathered}
\tag{PGGT.5}\label{eq:pggt-concrete-generator-control}
\]
Here every denominator is the actual spectral minimum of
the corresponding row of PGGT.2, not an additional
hypothesis. To prove the middle inequality, the absolutely
convergent finite-dimensional integral
\[
 Y=\int_0^\infty
 e^{-u\mathsf S}[B,T_t]e^{-u\mathsf S}\,du
\]
has $(i,j)$ block $(b_i-b_j)/(s_i+s_j)P_iT_tP_j$
and hence equals $[\mathsf S,T_t]$, since
$b_i-b_j=s_i^2-s_j^2$. The full repeated eigenspaces
remain present; their equal-eigenvalue blocks are zero.
Because $s_i\geq\sqrt\alpha$, the integral norm is
at most $\|[B,T_t]\|_G/(2\sqrt\alpha)$.
Multiplication by $\mathsf S^{-1}$ and adding its adjoint
give the middle inequality, including its factor one.
The first inequality is the reverse triangle inequality
after the exact isometry PGGT.3. For the last, put
$Q_j=\sum_{i>j}P_i$. Then
\[
 B=b_1I+\sum_{j=1}^{r-1}(b_{j+1}-b_j)Q_j.
\]
In the original $G$-orthogonal image/kernel decomposition
of $Q_j$, the commutator $[Q_j,T_t]$ has off-diagonal
blocks $T_{12}$ and $-T_{21}$ and zero diagonal blocks.
Its norm is $\max(\|T_{12}\|,\|T_{21}\|)\leq\|T_t\|_G$.
The triangle inequality and the telescoping sum of the
positive gaps prove the last bound. When $r=1$, the sum
is empty and every commutator and error above is zero.

The whole rank-one residue term remains explicit even
before taking its norm. Write $x=e_0$,
$y=G^{-1}\ell^*$; then $\ell(v)=\langle y,v\rangle_G$.
Self-adjointness of $B$ gives
\[
 [B,R]v=(Bx)\langle y,v\rangle_G
               -x\langle By,v\rangle_G,
 \qquad [B,R]=[Bx,x]\,[y,-By]^{\dagger_G}.
 \tag{PGGT.6}\label{eq:pggt-residue-sign}
\]
The displayed expression proves the factorization on each
original vector. Thus PGGT.4--5 contains its negative
term and its full interference with $[B,A]$ and the
original complex period parameter; no separate positive
bound replaces that actual sum.

For the completed row the $b_i$ in PGGT.3 are exactly the
distinct values among the full $\lambda_{k,j}$ in
PGMB.9. Thus its concrete last constant is
$(\lambda_{k,q_k}-\lambda_{k,1})/(1+\lambda_{k,1})$,
and its earlier whole-error and commutator controls retain
every eigenspace. At $k=1$, $B=0$, $\mathsf S=I$ and
the two Hermitian generators agree exactly. At $k\geq2$
strict positivity supplies the displayed positive
denominator; it supplies no unstated smallness of the
spectral diameter as $k$ varies.

Finally the completed generator has an exact realization
on the original Hilbert targets, rather than only on a
coefficient form. Retain
$\mathfrak U_k=(L_{\rm low},\mathscr D_k^{\rm mix})$
and $\mathfrak U_k^*\mathfrak U_k=\Omega_k$ from PGMB.6.
On its original direct-sum target define
\[
 \widetilde T_t=\mathfrak U_kT_t\Omega_k^{-1}\mathfrak U_k^*,
 \quad\widetilde T_t\mathfrak U_k=\mathfrak U_kT_t,
 \quad\ker\widetilde T_t
 =\mathfrak U_k(\ker T_t)\mathbin{\oplus^\perp}
                                  \mathfrak U_k(E)^\perp.
 \tag{PGGT.7}\label{eq:pggt-realized-generator}
\]
Every map is bounded because its intermediate coefficient
space is finite dimensional. The Gram identity proves
the intertwining. Every target vector decomposes uniquely
as $\mathfrak U_kx+v_\perp$, with $x=\Omega_k^{-1}
\mathfrak U_k^*v$; the formula then sends it to
$\mathfrak U_kT_tx$, proving the complete kernel.
Taking adjoints in the original direct sum gives
\[
 \widetilde T_t+\widetilde T_t^*
 =\mathfrak U_k(T_t+T_t^{\dagger_{\Omega_k}})
                                      \Omega_k^{-1}\mathfrak U_k^*,
 \qquad
 \|\widetilde T_t+\widetilde T_t^*\|
                           =\epsilon_{\Omega_k}(t).
 \tag{PGGT.8}\label{eq:pggt-realized-hermitian-control}
\]
Indeed $T_t^{\dagger_{\Omega_k}}=\Omega_k^{-1}T_t^*\Omega_k$,
so both sides expand to the same two ordered terms. The
isometry from $(E,\Omega_k)$ onto $\mathfrak U_k(E)$
identifies their restrictions, and both vanish on its
orthogonal complement. This proves equality of the
operator norms. PGGT.5 therefore controls this specified
realized Hermitian generator through the original
$K_{\rm low}$ generator with its exact mixed discrepancy.
The calculation uses the actually computed finite-dimensional
image. It asserts no identification with an unbounded
ambient multiplication operator outside that image.
